% ============================================================
% Section 67: 二元合金的混合熵、相互作用能与相分离
% ============================================================
\section{固体理论：二元合金的混合熵、相互作用能与相分离}
\label{sec:alloy-mixing-entropy}

\begin{modelbox}[题目：简立方二元合金的统计热力学与相分离]
考虑由 $N$ 个原子组成的简立方晶体，其中 $A$、$B$ 两种原子随机混合。$A$--$A$ 和 $B$--$B$ 近邻相互作用能为 $-J$，$A$--$B$ 近邻相互作用能为 $+J$（$J>0$）。设 $A$ 原子数为 $N_A=N(1+x)/2$，$B$ 原子数为 $N_B=N(1-x)/2$。
\begin{enumerate}
    \item 计算系统的位形熵 $S$；
    \item 用平均场近似计算系统的总相互作用能 $U$；
    \item 写出自由能 $F(x)$ 并求临界温度 $T_c$；
    \item 当 $T<T_c$ 时，求 Spinodal 曲线（失稳分解线）$x_p(T)$。
\end{enumerate}
\end{modelbox}

\begin{tipbox}[考点快速分析]
\begin{itemize}
    \item \textbf{位形熵}：混合熵 $S=k_B\ln W$，用斯特林近似化简
    \item \textbf{平均场近似}：任一格点周围出现 $A$ 或 $B$ 的概率等于各自的浓度
    \item \textbf{自由能展开}：小 $x$ 下对数项泰勒展开，得到朗道型自由能
    \item \textbf{临界温度}：$\partial^2 F/\partial x^2|_{x=0}=0$ 给出 $T_c$
    \item \textbf{Spinodal 分解}：$\partial^2 F/\partial x^2<0$ 的成分区间内，任何涨落均导致自发分相
\end{itemize}
\end{tipbox}

\begin{derivationbox}[标准解题过程]
\textbf{Step 1: 位形熵的计算}

$N$ 个格点中，$N_A$ 个为 $A$ 原子，$N_B$ 个为 $B$ 原子。排列方式数为
\begin{equation}
W = \binom{N}{N_A} = \frac{N!}{N_A!\,N_B!}.
\end{equation}
位形熵 $S=k_{\!B}\ln W$。利用斯特林近似 $\ln M!\approx M\ln M-M$：
\begin{align}
S &= k_{\!B}\bigl[N\ln N - N_A\ln N_A - N_B\ln N_B\bigr] \notag\\
&= k_{\!B}\Bigl[N\ln N - \frac{N(1+x)}{2}\ln\!\Bigl(\frac{N(1+x)}{2}\Bigr) - \frac{N(1-x)}{2}\ln\!\Bigl(\frac{N(1-x)}{2}\Bigr)\Bigr].
\end{align}
展开对数项 $\ln[N(1\pm x)/2]=\ln N + \ln(1\pm x)-\ln 2$，化简得
\begin{equation}
\boxed{S = Nk_{\!B}\ln 2 - \frac{Nk_{\!B}}{2}\Bigl[(1+x)\ln(1+x) + (1-x)\ln(1-x)\Bigr].}
\end{equation}

\textbf{Step 2: 总相互作用能的计算}

简立方结构中每个原子有 $6$ 个最近邻。在平均场近似下，任一格点周围出现 $A$ 原子的概率为 $N_A/N$，出现 $B$ 原子的概率为 $N_B/N$。

一个 $A$ 原子与周围 $6$ 个原子的平均相互作用能为
\begin{equation}
u_A = 6\Bigl[(-J)\frac{N_A}{N} + (+J)\frac{N_B}{N}\Bigr] = \frac{6J}{N}(-N_A+N_B).
\end{equation}
一个 $B$ 原子的平均相互作用能为
\begin{equation}
u_B = 6\Bigl[(-J)\frac{N_B}{N} + (+J)\frac{N_A}{N}\Bigr] = \frac{6J}{N}(-N_B+N_A).
\end{equation}

系统总相互作用能（每对键被计算两次，故乘 $1/2$）：
\begin{align}
U &= \frac{1}{2}(N_A\nu_A + N_B\nu_B)
= \frac{3J}{N}\bigl[-N_A^2 - N_B^2 + 2N_A N_B\bigr] \notag\\
&= -\frac{3J}{N}(N_A-N_B)^2.
\end{align}
代入 $N_A-N_B=Nx$：
\begin{equation}
\boxed{U = -3NJx^2.}
\end{equation}

\textbf{Step 3: 自由能展开与临界温度}

自由能 $F=U-TS$：
\begin{multline}
F(x) = -3NJx^2 - Nk_{\!B}T\ln 2 \\
+ \frac{Nk_{\!B}T}{2}\Bigl[(1+x)\ln(1+x) + (1-x)\ln(1-x)\Bigr].
\end{multline}

小 $x$ 下对数项泰勒展开：
\begin{equation}
(1+x)\ln(1+x) + (1-x)\ln(1-x) = x^2 + \frac{x^4}{6} + O(x^6).
\end{equation}
保留至 $x^4$ 项：
\begin{equation}
F(x) \approx -Nk_{\!B}T\ln 2 + \frac{1}{2}Nk_{\!B}(T-T_c)x^2 + \frac{1}{12}Nk_{\!B}Tx^4,
\end{equation}
其中
\begin{equation}
\boxed{T_c = \frac{6J}{k_{\!B}}.}
\end{equation}

\textit{稳定性分析}：
\begin{equation}
\frac{\partial^2 F}{\partial x^2}\bigg|_{x=0} = Nk_{\!B}(T-T_c).
\end{equation}
\begin{itemize}
    \item $T>T_c$：$\partial^2 F/\partial x^2|_{x=0}>0$，$x=0$ 为极小值，单相固溶体稳定
    \item $T<T_c$：$\partial^2 F/\partial x^2|_{x=0}<0$，$x=0$ 为极大值，系统不稳定，发生相分离
\end{itemize}

\textbf{Step 4: Spinodal 曲线与相成分}

当 $T<T_c$ 时，自由能曲线在 $x=0$ 两侧出现两个极小值点，对应平衡的两相成分。由自由能极小条件 $\partial F/\partial x=0$：
\begin{equation}
\frac{\partial F}{\partial x} = Nk_{\!B}(T-T_c)x + \frac{1}{3}Nk_{\!B}Tx^3 = 0.
\end{equation}
舍去 $x=0$ 的平庸解，得
\begin{equation}
(T-T_c) + \frac{T}{3}x^2 = 0
\quad\Longrightarrow\quad
\boxed{x_p(T) = \pm\sqrt{\frac{3(T_c-T)}{T}}.}
\end{equation}

\textit{物理意义}：在 $|x|<x_p(T)$ 的成分区间内，$\partial^2 F/\partial x^2<0$，系统处于热力学\textbf{不稳定状态}（Spinodal 区），任何微小的成分涨落都会导致系统自发分解为两相（成核生长或 Spinodal 分解）。
\end{derivationbox}

\begin{formulabox}[答案汇总]
\begin{center}
\begin{tabular}{ll}
\toprule
\textbf{问题} & \textbf{答案} \\
\midrule
(1) 位形熵 & $\displaystyle S = Nk_{\!B}\ln 2 - \frac{Nk_{\!B}}{2}\bigl[(1+x)\ln(1+x)+(1-x)\ln(1-x)\bigr]$ \\
(2) 总相互作用能 & $U = -3NJx^2$ \\
(3) 临界温度 & $T_c = 6J/k_{\!B}$ \\
(4) Spinodal 曲线 & $\displaystyle x_p(T) = \pm\sqrt{3(T_c-T)/T}$ \\
\bottomrule
\end{tabular}
\end{center}
\end{formulabox}

\begin{insightbox}[物理图像：为什么低温下会相分离？]
本题的核心竞争在于\textbf{能量}与\textbf{熵}的博弈：
\begin{itemize}
    \item \textbf{能量项}（$U=-3NJx^2$）：$J>0$ 意味着 $A$--$B$ 相互作用（$+J$）比 $A$--$A$、$B$--$B$（$-J$）更不利。系统倾向于\textbf{分离}为 $A$ 富集相和 $B$ 富集相，以降低能量
    \item \textbf{熵项}（混合熵）：系统倾向于\textbf{混合}，因为混合态的微观状态数更多
\end{itemize}

\textbf{温度决定胜负}：
\begin{itemize}
    \item 高温：熵主导，$F$ 在 $x=0$ 处取极小，\textbf{单相固溶体}稳定
    \item 低温：能量主导，$F$ 在 $x=0$ 处取极大，系统\textbf{自发分相}
    \item $T=T_c$：能量与熵势均力敌，临界点
\end{itemize}

\textbf{与 Ising 模型的联系}：本题本质上是 $z=6$（简立方配位数）的伊辛模型在平均场近似下的结果。临界温度 $T_c=zJ/k_{\!B}=6J/k_{\!B}$ 与平均场理论一致。
\end{insightbox}

\begin{thinkbox}[考场速记：二元合金相分离问题的通用框架]
遇到``混合熵''``临界温度''``Spinodal''类问题：
\begin{enumerate}
    \item \textbf{写位形熵}：$S=k_{\!B}\ln W = k_{\!B}[N\ln N - N_A\ln N_A - N_B\ln N_B]$
    \item \textbf{写相互作用能}：平均场近似下，$U\propto -(N_A-N_B)^2$ 或更一般地 $U\propto -x^2$
    \item \textbf{写自由能}：$F=U-TS$，对数项小 $x$ 展开：$(1\pm x)\ln(1\pm x)\approx x^2\pm x^3/3+x^4/6$
    \item \textbf{求临界温度}：$\partial^2 F/\partial x^2|_{x=0}=0$
    \item \textbf{求 Spinodal}：$\partial F/\partial x=0$ 的非零解
\end{enumerate}

\textbf{自由能曲线的形状判断}：
\begin{itemize}
    \item $T>T_c$：单谷，$x=0$ 为全局极小（混合相）
    \item $T<T_c$：双谷，$x=0$ 为局部极大（分相），两侧极小为两相平衡成分
\end{itemize}
\textit{陷阱提醒}：总相互作用能公式中的 $1/2$ 因子（每对键被计算两次）不要遗漏；泰勒展开时注意 $(1+x)\ln(1+x)$ 的展开式中 $x^3$ 项系数为 $+1/6$，而 $(1-x)\ln(1-x)$ 中为 $-1/6$，两者相加后奇次项抵消。
\end{thinkbox}
